\documentclass{article}%
\usepackage[T1]{fontenc}%
\usepackage[utf8]{inputenc}%
\usepackage{lmodern}%
\usepackage{textcomp}%
\usepackage{lastpage}%
\usepackage{geometry}%
\geometry{margin=1in,headheight=14pt,headsep=25pt}%
\usepackage{hyperref}%
\usepackage{enumitem}%
\usepackage{fancyhdr}%
\usepackage{xcolor}%
%

            \hypersetup{
                colorlinks=true,
                linkcolor=blue,
                filecolor=magenta,
                urlcolor=blue,
            }
            
            \pagestyle{fancy}
            \fancyhf{}
            \rhead{Generated on \today}
            \lhead{Course Summary}
            \cfoot{\thepage}
        %
\title{Course Summary}%
\author{Generated by ScribeLec}%
\date{\today}%
%
\begin{document}%
\normalsize%
\maketitle%
\tableofcontents\newpage%
\section{Key Terms}%
\label{sec:KeyTerms}%
\begin{itemize}[[leftmargin=*]]%
\item \href{https://www.math.purdue.edu/~yipn/421/2024-08-27-ExSimplex.pdf#page=1}{\textbf{Simplex Method}}: An iterative algorithm used to solve linear programming problems by moving from one vertex of the feasible region to another until an optimal solution is found.<L[2024-08-27-ExSimplex] 1><L[2024-08-27-ExSimplex] 6>%
\item \href{https://www.math.purdue.edu/~yipn/421/2024-08-27-ExSimplex.pdf#page=1}{\textbf{Slack Variable}}: A non-negative variable added to a less-than-or-equal-to constraint to transform it into an equality constraint.<L[2024-08-27-ExSimplex] 1>%
\item \href{https://www.math.purdue.edu/~yipn/421/2024-08-27-ExSimplex.pdf#page=2}{\textbf{Feasible Region}}: The set of all points that satisfy all the constraints of a linear programming problem.<L[2024-08-27-ExSimplex] 2><L[2024-08-27-ExSimplex] 4><L[2024-08-27-ExSimplex] 5><L[2024-08-27-ExSimplex] 30>%
\item \href{https://www.math.purdue.edu/~yipn/421/2024-08-27-ExSimplex.pdf#page=4}{\textbf{Vertex}}: A point where two or more constraint boundaries intersect in the feasible region of a linear programming problem. Vertices represent potential optimal solutions.<L[2024-08-27-ExSimplex] 4><L[2024-08-27-ExSimplex] 5><L[2024-08-27-ExSimplex] 30>%
\item \href{https://www.math.purdue.edu/~yipn/421/2024-08-27-ExSimplex.pdf#page=9}{\textbf{Basic Variable}}: In the Simplex method, a variable that is expressed in terms of non-basic variables in a dictionary.<L[2024-08-27-ExSimplex] 9><L[2024-08-27-ExSimplex] 10><L[2024-08-27-ExSimplex] 26><L[2024-08-27-ExSimplex] 27><L[2024-08-27-ExSimplex] 28>%
\item \href{https://www.math.purdue.edu/~yipn/421/2024-08-27-ExSimplex.pdf#page=6}{\textbf{Non{-}basic Variable}}: In the Simplex method, a variable set to zero in a given iteration of the algorithm.<L[2024-08-27-ExSimplex] 6><L[2024-08-27-ExSimplex] 7><L[2024-08-27-ExSimplex] 26><L[2024-08-27-ExSimplex] 27><L[2024-08-27-ExSimplex] 28>%
\item \href{https://www.math.purdue.edu/~yipn/421/2024-08-27-ExSimplex.pdf#page=9}{\textbf{Pivot Operation}}: The process of exchanging a basic and a non-basic variable in the Simplex method to move to a new vertex in the feasible region.<L[2024-08-27-ExSimplex] 9><L[2024-08-27-ExSimplex] 10><L[2024-08-27-ExSimplex] 12><L[2024-08-27-ExSimplex] 13><L[2024-08-27-ExSimplex] 14><L[2024-08-27-ExSimplex] 15>%
\item \href{https://www.math.purdue.edu/~yipn/421/2024-08-27-ExSimplex.pdf#page=26}{\textbf{Dictionary}}: A representation of a linear programming problem where basic variables are expressed as functions of non-basic variables.<L[2024-08-27-ExSimplex] 26><L[2024-08-27-ExSimplex] 27><L[2024-08-27-ExSimplex] 28><L[2024-08-27-ExSimplex] 29>%
\item \href{https://www.math.purdue.edu/~yipn/421/2024-08-27-ExSimplex.pdf#page=32}{\textbf{Auxiliary Problem}}: A modified linear programming problem introduced to find an initial feasible solution when the origin is not feasible in the original problem.<L[2024-08-27-ExSimplex] 32><L[2024-08-27-ExSimplex] 33><L[2024-08-27-ExSimplex] 34><L[2024-08-27-ExSimplex] 35><L[2024-08-27-ExSimplex] 36><L[2024-08-27-ExSimplex] 37><L[2024-08-27-ExSimplex] 38><L[2024-08-27-ExSimplex] 39><L[2024-08-27-ExSimplex] 40><L[2024-08-27-ExSimplex] 41><L[2024-08-27-ExSimplex] 42>%
\end{itemize}%
\newpage

%
\section{Problem Types}%
\label{sec:ProblemTypes}%
\begin{itemize}[[leftmargin=*]]%
\item \href{https://www.math.purdue.edu/~yipn/421/2024-08-27-ExSimplex.pdf#page=32}{\textbf{Finding an Initial Feasible Solution}}: The problem of finding a feasible solution when the origin is not in the feasible region is addressed by introducing an auxiliary problem. <L[2024-08-27-ExSimplex] 32><L[2024-08-27-ExSimplex] 33><L[2024-08-27-ExSimplex] 34><L[2024-08-27-ExSimplex] 35><L[2024-08-27-ExSimplex] 36><L[2024-08-27-ExSimplex] 37><L[2024-08-27-ExSimplex] 38><L[2024-08-27-ExSimplex] 39><L[2024-08-27-ExSimplex] 40><L[2024-08-27-ExSimplex] 41><L[2024-08-27-ExSimplex] 42>%
\item \href{https://www.math.purdue.edu/~yipn/421/2024-08-27-ExSimplex.pdf#page=2}{\textbf{Solving Linear Programs Graphically in Two Dimensions}}: This problem involves finding the optimal solution to a linear program by graphically identifying the feasible region and its vertices. <L[2024-08-27-ExSimplex] 2><L[2024-08-27-ExSimplex] 3><L[2024-08-27-ExSimplex] 4><L[2024-08-27-ExSimplex] 5><L[2024-08-27-ExSimplex] 16>%
\item \href{https://www.math.purdue.edu/~yipn/421/2024-08-27-ExSimplex.pdf#page=1}{\textbf{Solving Linear Programs Using the Simplex Method}}: This problem focuses on using the simplex method to iteratively improve a solution by moving from one vertex to another in the feasible region until the optimal solution is found. <L[2024-08-27-ExSimplex] 1><L[2024-08-27-ExSimplex] 6><L[2024-08-27-ExSimplex] 7><L[2024-08-27-ExSimplex] 8><L[2024-08-27-ExSimplex] 9><L[2024-08-27-ExSimplex] 10><L[2024-08-27-ExSimplex] 11><L[2024-08-27-ExSimplex] 12><L[2024-08-27-ExSimplex] 13><L[2024-08-27-ExSimplex] 14><L[2024-08-27-ExSimplex] 15><L[2024-08-27-ExSimplex] 17><L[2024-08-27-ExSimplex] 18><L[2024-08-27-ExSimplex] 19><L[2024-08-27-ExSimplex] 20><L[2024-08-27-ExSimplex] 21><L[2024-08-27-ExSimplex] 22><L[2024-08-27-ExSimplex] 23><L[2024-08-27-ExSimplex] 24><L[2024-08-27-ExSimplex] 25>%
\item \href{https://www.math.purdue.edu/~yipn/421/2024-08-27-ExSimplex.pdf#page=26}{\textbf{Using Dictionaries to Represent and Solve Linear Programs}}: This problem involves using dictionaries to represent the linear program, where basic variables are expressed in terms of non-basic variables, and iteratively updating the dictionary through pivot operations to find the optimal solution. <L[2024-08-27-ExSimplex] 26><L[2024-08-27-ExSimplex] 27><L[2024-08-27-ExSimplex] 28><L[2024-08-27-ExSimplex] 29>%
\item \href{https://www.math.purdue.edu/~yipn/421/2024-08-27-ExSimplex.pdf#page=30}{\textbf{Handling Infeasible Origins in Linear Programming}}: This problem addresses situations where the origin is not a feasible solution, requiring the use of an auxiliary problem to find an initial feasible solution before applying the simplex method. <L[2024-08-27-ExSimplex] 30><L[2024-08-27-ExSimplex] 31>%
\end{itemize}%
\newpage

%
\section{Algorithm Solutions}%
\label{sec:AlgorithmSolutions}%
\begin{itemize}[[leftmargin=*]]%
\item \href{https://www.math.purdue.edu/~yipn/421/2024-08-27-ExSimplex.pdf#page=6}{\textbf{Simplex Method}}: Iteratively move from one vertex of the feasible region to another by setting non-basic variables to zero and choosing a new set of (non)basic variables to improve the objective function. This corresponds to moving from vertex to vertex in the graphical representation.<L[2024-08-27-ExSimplex] 6><L[2024-08-27-ExSimplex] 7><L[2024-08-27-ExSimplex] 8><L[2024-08-27-ExSimplex] 9><L[2024-08-27-ExSimplex] 10><L[2024-08-27-ExSimplex] 11><L[2024-08-27-ExSimplex] 12><L[2024-08-27-ExSimplex] 13><L[2024-08-27-ExSimplex] 14><L[2024-08-27-ExSimplex] 15><L[2024-08-27-ExSimplex] 16><L[2024-08-27-ExSimplex] 17><L[2024-08-27-ExSimplex] 18><L[2024-08-27-ExSimplex] 19><L[2024-08-27-ExSimplex] 20><L[2024-08-27-ExSimplex] 21><L[2024-08-27-ExSimplex] 22><L[2024-08-27-ExSimplex] 23><L[2024-08-27-ExSimplex] 24><L[2024-08-27-ExSimplex] 25>%
\item \href{https://www.math.purdue.edu/~yipn/421/2024-08-27-ExSimplex.pdf#page=9}{\textbf{Pivot Operation}}: Interchange one basic and one non-basic variable in the dictionary to move to a new vertex in the feasible region, updating the equations and objective function accordingly.<L[2024-08-27-ExSimplex] 9><L[2024-08-27-ExSimplex] 10><L[2024-08-27-ExSimplex] 12><L[2024-08-27-ExSimplex] 13>%
\item \href{https://www.math.purdue.edu/~yipn/421/2024-08-27-ExSimplex.pdf#page=32}{\textbf{Auxiliary Problem}}: Introduce a new variable $x_0$ and modify constraints to create a new problem where the origin is feasible. Solve this problem to find a feasible solution for the original problem, then proceed with the Simplex method.<L[2024-08-27-ExSimplex] 32><L[2024-08-27-ExSimplex] 33><L[2024-08-27-ExSimplex] 34><L[2024-08-27-ExSimplex] 35><L[2024-08-27-ExSimplex] 36><L[2024-08-27-ExSimplex] 37><L[2024-08-27-ExSimplex] 38><L[2024-08-27-ExSimplex] 39><L[2024-08-27-ExSimplex] 40><L[2024-08-27-ExSimplex] 41><L[2024-08-27-ExSimplex] 42>%
\end{itemize}%
\newpage

%
\end{document}